Concluding 
The formalization of the non-stationary Susceptible-Infected-Mitigated (\(\mathcal{S}\mathcal{I}\mathcal{M}\)) model supplies a mathematically rigorous and structurally determinative account of multi-agent AI risk. By discarding anthropomorphic narratives of spontaneous code intent, the framework demonstrates that Compromised Swarms and Adversary-Adapted Swarms are dual manifestations of a single engineering failure: the unmitigated transmission of execution privileges across an inadequately guarded agent graph.

Summary of Core Structural Findings

Three mathematically validated axioms underwrite the thesis:

Topological Symmetries  
   Runtime hijacking and offline architectural exfiltration share identical supercritical transmission pathways. Cascade velocity is governed exclusively by the human-chosen, time-varying parameter set
   \[
   \bigl\{\langle k(t)\rangle,\;\langle T(t)\rangle,\;\bar\alpha(t),\;\gamma(t),\;\delta(t)\bigr\},
   \]
   not by any autonomous property of the underlying models.

Stochastic Vulnerabilities  
   Continuous fluid approximations break down under the extreme heterogeneity and finite-population noise of production clusters. Gillespie SSA realizations reveal a characteristic bimodal structure: an exploit either burns out among low-degree workers or executes a hub-amplification super-jump that collapses the entire topology in a handful of discrete events—even when global averages appear subcritical.

The Engineering Fix  
   Behavioral alignment, constitutional self-critique, and perimeter gateways leave the governing parameters untouched and therefore afford no reliable defense. The sole mathematically defensible remediation is continuous, graph-level enforcement that keeps the instantaneous basic reproduction number inside the hardened envelope
   \[
   R_0(t)\;\le\;0.5.
   \]

                      [ Multi-Agent System State Space ]
                                       │
                ┌──────────────────────┴──────────────────────┐
                ▼                                             ▼
  Supercritical Regime (R₀(t) > 1)              Subcritical Target (R₀(t) ≤ 0.5)
 ┌─────────────────────────────┐               ┌─────────────────────────────┐
 │  • Cascade proliferation    │               │  • Immediate decay          │
 │  • Coordinated tool chaining│               │  • Graph-node isolation     │
 │  • Hub super-jumps          │  ──[Fix]──►   │  • Deterministic safety     │
 │  • Structural isomorphism   │               │  • Cryptographic sanity     │
 └─────────────────────────────┘               └─────────────────────────────┘

In short, deviant swarm capabilities are not emergent properties of autonomous code; they are the macroscopic expression of architectural parameters that system designers themselves select and, too often, leave unmonitored. Restoring safety is therefore an engineering task: measure the full parameter vector in real time, actuate topological controls that force \(R_0(t)\) below the mandated ceiling, and thereby convert an inherently supercritical interaction graph into a provably subcritical one. When this discipline is observed, both the runtime and the offline vectors of the dual contagion are rendered mathematically impossible.
